% Project proposal for the professor. First draft written by Claude Code on 2026-10-10 at
% Laith's request, after the project moved from bank customer checks to visitor visa
% applications. Changed the same day after Laith's decisions: the applicants live in the
% United Arab Emirates; every application states the same legal status; the project is
% about bias among the 22 Arab League nationalities, measured against the official rules,
% and does not look at passport strength or at non-Arab comparison groups.
% Extended by Claude Code on 2026-10-12 at Laith's request: the paragraph on machines that
% sort visa applicants now quotes its sources word for word and adds Australia and the
% plans in the EU and Germany. The quotes were checked on the originals that day; how far
% each source was read is in research/literature_matrix.csv.
% Read it, correct it and put it into your own words before sending. The
% proposal of 2026-10-08 (bank customer checks) is in
% research/old_design/bank_2026-10-08/proposal/; the one of 2026-10-04 (hiring) is in
% research/old_design/hiring_2026-10-04/proposal/.
\documentclass[11pt]{article}

\usepackage[preprint]{acl}

\usepackage{times}
\usepackage{latexsym}
\usepackage[T1]{fontenc}
\usepackage[utf8]{inputenc}
\usepackage{microtype}
\usepackage{inconsolata}
\usepackage{graphicx}
\usepackage{booktabs}
\usepackage{enumitem}
\setlist{itemsep=1pt,topsep=3pt,parsep=0pt}

\title{Which Arab? How LLMs Rate Visa Applicants from the 22 Arab League Countries\\[4pt]
\large Project proposal, IS 617 (HWS26)}

\author{
  Laith Sandouk \and Aleksander Kasak \and Nico Fedotov \\
  University of Mannheim \\
  \texttt{\{laith.sandouk, aleksander.kasak, nico.fedotov\}@students.uni-mannheim.de}
}

\begin{document}
\maketitle

\section{Research questions}
\label{sec:rq}

Our motivation is Islamophobia. Muslims in Europe report discrimination in growing
numbers, and language models have been shown to repeat stereotypes about Muslims
(Section~\ref{sec:related}). We want to know whether this shows when a model takes part
in a decision about a person. A decision of that kind rarely has a religion in front of
it. At a border it has a passport. All Muslim-majority countries together are too mixed
a group for one experiment: 53 countries on three continents
\citep{pew2025religious}. We therefore study one group among them that shares a language
and a label, the 22 members of the Arab League, and we ask about bias among them. Two
people hand in the same application for a visitor visa. One holds a Syrian passport, the
other an Emirati one. The official rules for the visa are the same for both. Why should
a model rate them differently? And if it does, does it sort the 22 countries into the
groups a reader would expect: rich, normal and poor? In the experiment a model receives
the guidance of a visa office and one application. The applications are identical except
for the nationality and the passport, and the model returns a risk rating and a
recommendation.

\begin{description}[leftmargin=2.6em,labelwidth=2.1em,labelsep=0.5em]
  \item[RQ1] Do risk ratings and recommendations differ across the 22 Arab League
        nationalities when the application is identical and the official rules are the
        same for all of them? This is the first main question.
  \item[RQ2] Can the 22 countries be grouped into rich, normal and poor by how the models
        rate them? This is the second main question.
  \item[RQ3] Do the differences remain when the guidance states that nationality is not
        a ground for a higher rating?
  \item[RQ4] What would the differences cost an applicant: in the savings needed to make
        up for them, in fees and in extra checks?
\end{description}

\noindent Every question is answered for four models.

\section{Motivation and related work}
\label{sec:related}

\textbf{Anti-Muslim discrimination is widespread, and models repeat it.} In a survey in
13 EU countries, 47\% of Muslim respondents reported racial discrimination in the five
years before the survey, against 39\% in 2016; in Germany it was 68\%
\citep{fra2024being}. A German monitoring network documented 3,080 anti-Muslim incidents
in 2024 and 1,926 in 2023; part of the rise comes from a larger number of reporting
centres \citep{claim2025lagebild}. In GPT-3, 66 of 100 completions of the prompt ``Two
Muslims walked into a'' were violent \citep{abid2021persistent}, and a weaker form of
this association was still found in later GPT models \citep{hemmatian2023muslim}.

\textbf{From a religion to a passport.} In a hiring experiment in five European
countries, applicants whose families came from Muslim-majority countries were at a
disadvantage in four of the five even when they gave no sign of their religion
\citep{distasio2021muslim}. This is the link we rely on: a nationality is a noisy signal
for a religion. It is not the same thing. In a survey experiment in the United States,
respondents penalised the label ``Muslim'' but not the label ``Arab''
\citep{calfano2021evaluating}. Our applications state no religion, so we measure how
passports are treated.

\textbf{Our yardstick is the official rules.} We do not ask how strong a passport is or
how a country is seen. We ask what the rules of the destination say about the 22
nationalities. For the Australian Visitor visa they say the same for all of them. None
of the 22 can use the two simpler electronic permits
\citep{homeaffairs2026eta,homeaffairs2026evisitor}, so all need the same visa, apply
with the same online form \citep{au2021lin21056} and pay the same fee. The legal test is
about the individual and is the same for every nationality: the applicant must genuinely
intend to stay temporarily and must have adequate means \citep{au1994migrationregs}. The
law allows a list of risk groups defined by nationality, but the last such list was
revoked in 2009 \citep{au2008riskfactor,au2009revocation}. The department's instruction
for officers calls a profile ``merely an alert'', says that it ``is not grounds to refuse
to grant a visa'' and that ``all applications must be considered on their own merits''
\citep{homeaffairs2018genguide}; the version we read is from 2016, the current one is
not public. By these rules two identical applications should receive the same
assessment whatever the passport. A systematic difference among the 22 is what we call
bias.

\textbf{Why Australia.} It is the only destination we checked where the rules are the
same for all 22. Elsewhere the rules themselves sort the nationalities: Emirati citizens
need no visa, or only an electronic permit, for the Schengen area, the United Kingdom,
Canada and New Zealand, and the United States bars six of the 22 nationalities by
proclamation \citep{crs2025proclamation}. There a different answer for a Syrian and an
Emirati applicant follows from a rule and is not a bias.

\textbf{Same rules, different outcomes.} In practice the 22 fare very differently.
Australia granted 96.5\% of the visitor visa applications of Omani citizens and 21.9\%
of those of Somali citizens in 2019--20, the last year for which figures by citizenship
are available \citep{homeaffairs2020foi106}. Granted applications from the Gulf states
were decided in about one day at the median, Syrian and Palestinian ones in about a
month \citep{homeaffairs2025da449}. These figures come from real applicants in very
different circumstances, so they do not show a bias by themselves. They show what a
model may have learned. The experiment holds the applicant fixed and asks whether a
model carries the gap over to identical applications.

\textbf{Machines already sort visa applicants, and nationality has been an input.} We
quote the sources here, because the wording matters. From 2015 to 2020 a tool of the
United Kingdom streamed visa applications ``\,`Green' (low risk), `Amber' (medium risk)
or `Red' (high risk)''. In one office in early 2017, 96\% of the visit applications rated
green were issued and 49\% of those rated red, and the inspector saw ``a risk that the
streaming tool becomes a de facto decision-making tool'' \citep{icibi2017entry}. A
government letter quoted in a legal blog confirmed that ``nationality is one of the
relevant factors used by the streaming tool''; the tool was withdrawn in 2020 after a
legal challenge \citep{ukhrb2020streaming}. A Dutch tool builds profiles from
nationality, age group and gender, among other things \citep{nlminister2023iob}: ``33
percent of applications in the intensive, `high risk' track were rejected, whereas only
3.5 percent of applications in the normal, `low risk' track were rejected''
\citep{lighthouse2023ethnic}. Canada's system ``sorts applications into tiers based on
their level of complexity'', for all visitor visa applications made from abroad, and
approves routine ones; it ``never refuses applications nor does it recommend refusals''
\citep{ircc2022aia}. Australia, the destination in our experiment, uses analytics based
on machine learning to ``predict risk in visa programs'' and states that it ``does not
use machine learning to make administrative decisions'' \citep{homeaffairs2024foi}; the
document does not say what goes into the prediction.

\textbf{More is planned, and none of it is a language model.} EU law defines the
screening rules of its travel authorisation for visitors without a visa as ``an
algorithm enabling profiling'', built on data such as ``age range, sex, nationality''
\citep{eu2018etias}; the system was scheduled to start in the last quarter of 2026
\citep{meier2026algorithmic}. In Germany, a government paper calls the use of
artificial intelligence in the visa procedure ``den ersten Schritt'', the first step
\citep{aabmi2025kieckpunkte}, and a bill on artificial intelligence in migration
offices had its first reading on 8 October 2026 \citep{bundestag2026kimvg}. The British
tool was a decision tree and the Dutch one is, in the minister's words, not
self-learning; only the Australian document names machine learning. We found no
authority that lets a language model rate applications. Canada states that it ``does
not allow the use of generative AI tools, such as ChatGPT, in support of decision-making
on client applications'' \citep{ircc2025qpnote}. The EU AI Act classes systems that
assist in examining visa applications as high-risk \citep{eu2024aiact}. We do not test a
product; we test what open models do in the seat these systems occupy.

\textbf{What is known about language models.} In the closest study, GPT-3.5 and GPT-4
chose between pairs of immigrants from ten countries, three of them Arab League members;
most country effects were not significant, yet the models suggested additional scrutiny
for applicants from Iraq, Sudan and Somalia \citep{mao2025gatekeepers}. A public test
set for discrimination in model decisions contains a visitor visa case but varies age,
gender and race, not nationality
\citep{tamkin2023evaluating,anthropic2023discrimeval}. In a mortgage experiment with one
line stating the applicant's race, GPT-4 Turbo recommended approval 8.5 percentage
points less often for Black applicants, an instruction to use no bias removed that gap,
and the authors expressed the gap in credit-score points \citep{bowen2025measuring}.
Explicit labels can also have the opposite effect \citep{bunel2026customer}. We have
found no study that changes only the nationality in a visa application, and none that
compares the 22 nationalities. We have not searched systematically and treat this as
tentative.

\section{Design}
\label{sec:design}

\textbf{Materials.} Everything is synthetic; the project audits model behaviour and is
not a tool for assessing applicants. We prepared 20 applications for an Australian
Visitor visa (subclass 600, Tourist stream), with occupations from warehouse associate
to software developer. Ten applications contain nothing unusual. Ten contain one mild
point that the guidance names, such as a new job or a sibling who lives in Australia, so
that the rating has room to move. An application carries a reference number instead of a
name.

\textbf{Where the applicants live.} Every applicant was born in the United Arab
Emirates, has lived there since birth, is employed there, plans a holiday and has no
refusal or overstay on record. We chose the Emirates as the most realistic place for
people of all these nationalities to be born and raised without becoming citizens: birth
there gives no citizenship to the child of foreign parents \citep{uae1972nationality},
and 88.5\% of the residents were foreigners in the last official figures that give the
split \citep{khaleej2011expats}. Every application states the same legal status: a
lawful long-term resident with the full right to live and work there. Such a status
exists for foreigners, for example a renewable visa of five or ten years that needs no
sponsor \citep{uae2026goldenvisa}. The line simplifies reality for six nationalities: an
Emirati is a citizen there, and citizens of the five other Gulf states need no residence
visa \citep{uae2026gcc}. A model may know this although the application does not say
it, so we repeat the comparisons without these six as a check.

\textbf{Manipulation.} The nationality and the matching passport, for example
\texttt{Nationality: Jordanian} and \texttt{Jordanian passport}. Nothing else changes.
There are 26 conditions: the 22 Arab League nationalities, and four unrelated
nationalities (Uruguayan, Malawian, Maldivian, Cambodian) that need the same visa under
the same rules. The four are a noise check: any set of labels moves a model's answers a
little, and they show how much. They are not a research question.

\textbf{Three groups (RQ2).} Before any run we sort the 22 countries into three groups
by the World Bank's income classification \citep{worldbank2026wdi}. Rich (high income):
Bahrain, Kuwait, Oman, Qatar, Saudi Arabia and the United Arab Emirates. Normal (upper
and lower middle income): Algeria, Comoros, Djibouti, Egypt, Iraq, Jordan, Lebanon,
Libya, Mauritania, Morocco, Palestine and Tunisia. Poor (low income): Somalia, Sudan,
Syria and Yemen. Two limits are built in. The rich group is exactly the six Gulf
states, whose citizens have a stronger position in the Emirates in reality, so we cannot
separate ``rich'' from ``Gulf''. And the four poor countries are all at war or just out
of one.

\textbf{Guidance.} The model is told that it is a case assistant who pre-assesses
applications and that a visa officer takes every decision; this is the role the real
systems have. It receives a one-page guidance that paraphrases the public criteria. In a
second version the guidance has one more paragraph, built from the department's own
wording: nationality is not a ground for a higher rating, and every application must be
considered on its own merits (RQ3).

\textbf{Task.} For each application the model returns a risk score from 0 to 100, a
rating (low, medium, high), a list of further checks chosen from a menu of six, and a
recommendation (grant, refer for further checks, refuse). The request is phrased in three
ways.

\textbf{Cost (RQ4).} Our main measure is the amount of savings that would make up for a
passport: how much more money in the bank an applicant with passport~A needs to receive
the same risk score as the same applicant with passport~B. For this the applications
come with five levels of savings, and one block of runs repeats every application at
each level. Two further measures are the fees lost per 1,000 identical applicants (the
fee of AUD~250 is not refunded after a refusal) and the number of further checks.

\textbf{Models.} Four open-weight models served with vLLM on bwUniCluster: Qwen3 8B,
Llama 3.1 8B, Gemma 3 12B and Mistral Small 3.2 24B.

\textbf{Size.} Table~\ref{tab:runs} lists the runs per model.

\begin{table}[t]
\centering
\small
\begin{tabular}{@{}lr@{}}
\toprule
Block & Runs per model \\
\midrule
1\quad Guidance, no paragraph on nationality & 7,800 \\
2\quad Guidance with that paragraph & 7,800 \\
3\quad Block 1, four other savings levels & 6,240 \\
\midrule
Total & 21,840 \\
\bottomrule
\end{tabular}
\caption{Planned runs per model. Blocks 1 and 2: 20 applications, 26 nationalities, three
phrasings and five repetitions. Block 3: four further savings levels, three phrasings,
one run each.}
\label{tab:runs}
\end{table}

\section{What we will compare}
\label{sec:analysis}

\begin{description}[leftmargin=2.6em,labelwidth=2.1em,labelsep=0.5em]
  \item[RQ1] How far the risk score, the ratings, the further checks and the
        recommendations differ across the 22 Arab League nationalities. By the rules the
        difference should be zero. The four unrelated nationalities show how much a set
        of labels moves the answers by chance.
  \item[RQ2] Two checks. Do the three groups differ, and in the order rich, normal,
        poor? And the other way round: if the 22 are sorted into three groups by the
        models' ratings alone, do these match the three wealth groups, and which
        countries land in another group? Both for each model, phrasing and application.
  \item[RQ3] The same comparisons with the added paragraph.
  \item[RQ4] The three cost measures.
\end{description}

\noindent We have not fixed the statistical analysis yet. We will write down the
hypotheses and the analysis before the first real run.

\section{Timeline}
\label{sec:timeline}

Pitch on 13 October. A small pilot and block~1 on all four models before the midterm
presentation (11--12 November). Blocks~2 and~3 before the final presentation (9--10
December). Report due on 8 January 2027.

\section{Questions}
\label{sec:questions}

\begin{enumerate}
  \item Is the pre-assessment of visa applications a suitable setting, although no
        authority is documented as using a language model for it?
  \item Our motivation is Islamophobia, and what we vary is the passport of an Arab
        League country. Is this step convincing, or does the design need a condition
        that states a religion?
  \item Our yardstick is the official rules: they are the same for all 22
        nationalities, so we call any systematic difference between identical
        applications a bias. Is this convincing, given that real offices grant these
        nationalities at very different rates?
  \item We set the applicants in the United Arab Emirates and state the same legal
        status for everybody, although Gulf citizens have a stronger position there in
        reality. Is this setting convincing?
  \item We sort the 22 into rich, normal and poor by the World Bank's income classes (6,
        12 and 4 countries). Is this the right grouping, or should upper and lower
        middle income stay apart? Are four poor countries enough?
  \item Are four unrelated nationalities enough as a noise check, or is the check not
        needed at all?
  \item Is it acceptable to express a difference as the savings needed to make up for
        it, and is the extra block worth its runs?
  \item Are 20 applications, five repetitions and four open-weight models between 8B and
        24B parameters enough?
  \item Do you see concerns with reporting results by individual nationality?
\end{enumerate}

\paragraph{AI use.} The applications, the guidance text and the prompts, the table of
official rules, the literature and source search, and a first draft of this proposal
were prepared with the help of Claude Code.

\bibliography{../final_paper/custom}

\end{document}
